\documentclass[11pt,letterpaper]{article}
\usepackage{cornell-notes}
\usepackage{needspace}

\newcommand{\CornellDocumentTitle}{Chapter 99: Verifiable Voting Systems}
\title{\CornellDocumentTitle}
\author{}
\date{}
\setDocTitle{\CornellDocumentTitle}
\setDocSubtitle{Cybersecurity Cornell Notes}
\setCornellCollection{Cybersecurity Reference Notes}
\setCornellUnitType{Chapter}
\setCornellUnitNumber{99}
\setCornellUnitTitle{Verifiable Voting Systems}
\setCornellCompanionLabel{Cybersecurity study companion}

\begin{document}
\maketitle
\begin{CornellOverviewBox}{Chapter Orientation}
\textbf{Chapter orientation.} This chapter examines verifiable voting systems within the broader domain of advanced security. Its central problem is how to reason about voter, voting, ballot, and schemes while keeping security objectives, operating assumptions, evidence, and recovery connected. A practitioner should approach the material through physical consequences, safety, availability, environmental dependencies, remote connectivity, maintenance, and recovery. Useful prerequisites include basic risk, threat, control, and systems concepts.
\end{CornellOverviewBox}
\section*{Learning Objectives}
\begin{itemize}
\item Explain the security purpose and scope of Verifiable Voting Systems.
\item Identify the assets, actors, assumptions, and trust boundaries associated with voter, voting, and ballot.
\item Distinguish Introduction from Security Requirements and explain where their responsibilities intersect.
\item Analyze how failures in voter, voting, and ballot can affect confidentiality, integrity, availability, privacy, or safety.
\item Evaluate relevant controls through physical consequences, safety, availability, environmental dependencies, remote connectivity, maintenance, and recovery.
\item Audit an implementation using asset and dependency maps, physical-access records, sensor telemetry, safety constraints, maintenance logs, and recovery exercises.
\item Prioritize remediation according to exploitability, consequence, control strength, and recovery readiness.
\item Verify that corrective actions change observable risk rather than documentation alone.
\end{itemize}
\section*{Cornell Cue-and-Notes}
\Needspace{8\baselineskip}
\subsection*{Introduction}
\begin{CornellNotesTable}
\CornellNoteRow{What security problem does Introduction address?}{\textbf{Core idea.} Treat Introduction as a structured part of Verifiable Voting Systems, with particular attention to secret, vote, tion, and schemes. \textbf{Analytical lens.} This topic establishes a component of the chapter's argument; connect its definitions and mechanisms to a testable security objective. For secret, vote, and tion, model safety and physical consequences alongside conventional information-security effects. \textbf{Security reasoning.} Identify what must remain protected, who or what can alter the expected state, which assumptions cross a trust boundary, and how failure changes confidentiality, integrity, availability, privacy, or safety. \textbf{Application.} The practitioner should evaluate cyber and physical failure together, enforce safe states, and test operational recovery. \textbf{Evidence.} Seek asset and dependency maps, physical-access records, sensor telemetry, safety constraints, maintenance logs, and recovery exercises.}
\end{CornellNotesTable}
\begin{CornellSummaryBox}{Topic Summary}
Introduction is best remembered as the connection between secret, vote, tion, and schemes and the chapter's larger control objective. A defensible review states the assumption, expected behavior, evidence, failure consequence, and required response.
\end{CornellSummaryBox}
\Needspace{8\baselineskip}
\subsection*{Security Requirements}
\begin{CornellNotesTable}
\CornellNoteRow{Which assumptions matter in Security Requirements?}{\textbf{Core idea.} Treat Security Requirements as a structured part of Verifiable Voting Systems, with particular attention to voting, vote, encrypted, and requirements. \textbf{Analytical lens.} This topic establishes a component of the chapter's argument; connect its definitions and mechanisms to a testable security objective. For voting, vote, and encrypted, model safety and physical consequences alongside conventional information-security effects. \textbf{Security reasoning.} Identify what must remain protected, who or what can alter the expected state, which assumptions cross a trust boundary, and how failure changes confidentiality, integrity, availability, privacy, or safety. \textbf{Application.} The practitioner should evaluate cyber and physical failure together, enforce safe states, and test operational recovery. \textbf{Evidence.} Seek asset and dependency maps, physical-access records, sensor telemetry, safety constraints, maintenance logs, and recovery exercises. Related subtopics include Interrelationships and Conflicts, Challenges, Compromises, and Achieving System Security.}
\end{CornellNotesTable}
\begin{CornellSummaryBox}{Topic Summary}
Security Requirements is best remembered as the connection between voting, vote, encrypted, and requirements and the chapter's larger control objective. A defensible review states the assumption, expected behavior, evidence, failure consequence, and required response.
\end{CornellSummaryBox}
\Needspace{8\baselineskip}
\subsection*{Verifiable Voting Schemes}
\begin{CornellNotesTable}
\CornellNoteRow{How should a reviewer evaluate Verifiable Voting Schemes?}{\textbf{Core idea.} Treat Verifiable Voting Schemes as a structured part of Verifiable Voting Systems, with particular attention to vote, encrypted, fortune, and voting. \textbf{Analytical lens.} This topic establishes a component of the chapter's argument; connect its definitions and mechanisms to a testable security objective. For vote, encrypted, and fortune, model safety and physical consequences alongside conventional information-security effects. \textbf{Security reasoning.} Identify what must remain protected, who or what can alter the expected state, which assumptions cross a trust boundary, and how failure changes confidentiality, integrity, availability, privacy, or safety. \textbf{Application.} The practitioner should evaluate cyber and physical failure together, enforce safe states, and test operational recovery. \textbf{Evidence.} Seek asset and dependency maps, physical-access records, sensor telemetry, safety constraints, maintenance logs, and recovery exercises. Related subtopics include Verifiable Supervised Schemes.}
\end{CornellNotesTable}
\begin{CornellSummaryBox}{Topic Summary}
Verifiable Voting Schemes is best remembered as the connection between vote, encrypted, fortune, and voting and the chapter's larger control objective. A defensible review states the assumption, expected behavior, evidence, failure consequence, and required response.
\end{CornellSummaryBox}
\Needspace{8\baselineskip}
\subsection*{Building Blocks}
\begin{CornellNotesTable}
\CornellNoteRow{Why can Building Blocks fail in practice?}{\textbf{Core idea.} Treat Building Blocks as a structured part of Verifiable Voting Systems, with particular attention to key, secret, proof, and elgamal. \textbf{Analytical lens.} This topic establishes a component of the chapter's argument; connect its definitions and mechanisms to a testable security objective. For key, secret, and proof, model safety and physical consequences alongside conventional information-security effects. \textbf{Security reasoning.} Identify what must remain protected, who or what can alter the expected state, which assumptions cross a trust boundary, and how failure changes confidentiality, integrity, availability, privacy, or safety. \textbf{Application.} The practitioner should evaluate cyber and physical failure together, enforce safe states, and test operational recovery. \textbf{Evidence.} Seek asset and dependency maps, physical-access records, sensor telemetry, safety constraints, maintenance logs, and recovery exercises. Related subtopics include Verifiable Remote Schemes, Encryption Schemes, Rivest, Shamir, and Adelman Cipher, and ElGamal Cipher.}
\end{CornellNotesTable}
\begin{CornellSummaryBox}{Topic Summary}
Building Blocks is best remembered as the connection between key, secret, proof, and elgamal and the chapter's larger control objective. A defensible review states the assumption, expected behavior, evidence, failure consequence, and required response.
\end{CornellSummaryBox}
\Needspace{8\baselineskip}
\subsection*{Survey Of Noteworthy Schemes}
\begin{CornellNotesTable}
\CornellNoteRow{What security problem does Survey Of Noteworthy Schemes address?}{\textbf{Core idea.} Treat Survey Of Noteworthy Schemes as a structured part of Verifiable Voting Systems, with particular attention to voter, ballot, voting, and schemes. \textbf{Analytical lens.} This topic establishes a component of the chapter's argument; connect its definitions and mechanisms to a testable security objective. For voter, ballot, and voting, model safety and physical consequences alongside conventional information-security effects. \textbf{Security reasoning.} Identify what must remain protected, who or what can alter the expected state, which assumptions cross a trust boundary, and how failure changes confidentiality, integrity, availability, privacy, or safety. \textbf{Application.} The practitioner should evaluate cyber and physical failure together, enforce safe states, and test operational recovery. \textbf{Evidence.} Seek asset and dependency maps, physical-access records, sensor telemetry, safety constraints, maintenance logs, and recovery exercises. Related subtopics include Schemes Based on Mixnets, Schemes Based on Blind Signature, Specific Voter-Verifiable Schemes, and Schemes Based on Homomorphic Encryption.}
\end{CornellNotesTable}
\begin{CornellSummaryBox}{Topic Summary}
Survey Of Noteworthy Schemes is best remembered as the connection between voter, ballot, voting, and schemes and the chapter's larger control objective. A defensible review states the assumption, expected behavior, evidence, failure consequence, and required response.
\end{CornellSummaryBox}
\Needspace{8\baselineskip}
\subsection*{Threats To Verifiable Voting Systems}
\begin{CornellNotesTable}
\CornellNoteRow{Which assumptions matter in Threats To Verifiable Voting Systems?}{\textbf{Core idea.} Treat Threats To Verifiable Voting Systems as a structured part of Verifiable Voting Systems, with particular attention to ballot, voter, attack, and receipts. \textbf{Analytical lens.} This topic is threat-centered: separate actor capability, reachable attack surface, prerequisites, execution path, and consequence. For ballot, voter, and attack, model safety and physical consequences alongside conventional information-security effects. \textbf{Security reasoning.} Identify what must remain protected, who or what can alter the expected state, which assumptions cross a trust boundary, and how failure changes confidentiality, integrity, availability, privacy, or safety. \textbf{Application.} The practitioner should evaluate cyber and physical failure together, enforce safe states, and test operational recovery. \textbf{Evidence.} Seek asset and dependency maps, physical-access records, sensor telemetry, safety constraints, maintenance logs, and recovery exercises. Related subtopics include Coercion, Authentication of Receipts, Chain-Voting, and Use of Cryptography.}
\end{CornellNotesTable}
\begin{CornellSummaryBox}{Topic Summary}
Threats To Verifiable Voting Systems is best remembered as the connection between ballot, voter, attack, and receipts and the chapter's larger control objective. A defensible review states the assumption, expected behavior, evidence, failure consequence, and required response.
\end{CornellSummaryBox}
\begin{CornellWarningBox}{Historical Context}
Some products, protocols, examples, threat observations, or legal references in this chapter are historically bounded. Retain the underlying threat model and control objective, but verify present applicability before treating a named implementation as current guidance.
\end{CornellWarningBox}
\begin{CornellOverviewBox}{Defensive Guidance}
Apply the chapter by making assets, authority, trust, expected behavior, and failure response explicit. In practice, evaluate cyber and physical failure together, enforce safe states, and test operational recovery. A control is not fully supported until its operation, monitoring, ownership, and recovery behavior are demonstrated with evidence.
\end{CornellOverviewBox}
\section*{Threat-and-Control Map}
\begingroup\scriptsize\setlength{\tabcolsep}{2pt}
\begin{longtable}{>{\raggedright\arraybackslash}p{0.135\textwidth}>{\raggedright\arraybackslash}p{0.145\textwidth}>{\raggedright\arraybackslash}p{0.155\textwidth}>{\raggedright\arraybackslash}p{0.145\textwidth}>{\raggedright\arraybackslash}p{0.16\textwidth}>{\raggedright\arraybackslash}p{0.17\textwidth}}
\toprule\rowcolor{Navy}\color{white}\textbf{Asset} & \color{white}\textbf{Threat / Failure} & \color{white}\textbf{Vulnerability} & \color{white}\textbf{Impact} & \color{white}\textbf{Detection / Evidence} & \color{white}\textbf{Control}\\\midrule
\endfirsthead
\toprule\rowcolor{Navy}\color{white}\textbf{Asset} & \color{white}\textbf{Threat / Failure} & \color{white}\textbf{Vulnerability} & \color{white}\textbf{Impact} & \color{white}\textbf{Detection / Evidence} & \color{white}\textbf{Control}\\\midrule
\endhead
People, physical processes, connected devices, and essential services & Tampering, unsafe command, service disruption, surveillance, or cascading failure & Insecure remote access, fragile dependencies, weak update paths, or insufficient safety separation & Safety events, prolonged outage, physical damage, privacy harm, or public disruption & Operational telemetry, integrity monitoring, physical controls, and anomaly investigation & Layered cyber-physical protection, safe-state design, segmentation, secure maintenance, and rehearsed recovery\\\midrule
Assumptions surrounding voter & Misuse or unexpected state transition & Trust in voting is not validated & A local weakness becomes a broader compromise path & Review evidence for ballot and test negative paths & Make trust explicit, constrain authority, and fail safely\\\midrule
Recovery capability and decision evidence & Control failure remains unnoticed or cannot be reversed & Monitoring, ownership, or restoration has not been exercised & Longer exposure and slower, less reliable recovery & Alert-to-response exercises and restoration tests & Assigned ownership, actionable telemetry, preserved evidence, and rehearsed recovery\\\midrule
\bottomrule
\end{longtable}
\endgroup
\section*{Practical Security Checklist}
\begin{CornellChecklist}
\item Define the in-scope assets, users, services, data, and operational dependencies for Verifiable Voting Systems.
\item Document the expected behavior and security objective of voter.
\item Map identities, privileges, trust boundaries, and externally controlled inputs associated with voter and voting.
\item Record the threat or failure being addressed, its prerequisites, likely path, and credible consequences.
\item Inspect asset and dependency maps, physical-access records, sensor telemetry, safety constraints, maintenance logs, and recovery exercises.
\item Test both the intended path and negative paths, including invalid state, partial failure, excessive privilege, and unavailable dependencies.
\item Confirm preventive controls are complemented by actionable detection and a named response owner.
\item Exercise containment, restoration, and ballot; record actual recovery behavior and unresolved dependencies.
\item Validate exceptions, compensating controls, expiration dates, and accountable approval.
\item Preserve reproducible evidence and tie each finding to affected assets, risk, remediation, and verification criteria.
\end{CornellChecklist}
\begin{CornellSummaryBox}{Chapter Synthesis}
The chapter's principal lesson is that verifiable voting systems must be evaluated as an operating security system rather than an isolated product or statement. The most important failure patterns involve voter, voting, ballot, and schemes, especially when ownership, boundaries, telemetry, or recovery remain implicit. A reviewer should connect architecture and behavior to asset and dependency maps, physical-access records, sensor telemetry, safety constraints, maintenance logs, and recovery exercises, prioritize findings by reachable consequence and control independence, and verify remediation through observable change. Within Advanced Security, this chapter contributes a reusable method: state the objective, expose assumptions, test failure, preserve evidence, and learn from operation.
\end{CornellSummaryBox}
\section*{Self-Test Questions}
\begin{CornellRecallList}
\item What is the central security objective of Verifiable Voting Systems?
\item How does Introduction contribute to the chapter's broader security argument?
\item Distinguish Security Requirements from Verifiable Voting Schemes.
\item Which trust assumptions should be made explicit when evaluating voter?
\item How could a weakness involving voting affect confidentiality, integrity, availability, privacy, or safety?
\item What evidence would demonstrate that controls surrounding ballot operate as intended?
\item Why should prevention, detection, response, and recovery be evaluated together in this domain?
\item Scenario: A connected physical process remains functionally available after a cyber event but no longer operates within safe constraints. What should the reviewer examine first, and why?
\item Scenario: A control associated with schemes passes a configuration check but fails during an operational exercise. How should the finding be interpreted and prioritized?
\item How should a practitioner distinguish enduring principles in this chapter from time-sensitive tools, examples, or implementation details?
\end{CornellRecallList}
\Needspace{8\baselineskip}
\section*{Answer Key}
\begin{enumerate}
\item The objective is to protect the relevant assets by understanding physical consequences, safety, availability, environmental dependencies, remote connectivity, maintenance, and recovery and by connecting controls to measurable risk reduction.
\item Introduction provides one part of the causal chain: it should identify expected behavior, dependencies, failure modes, and the evidence needed to justify confidence.
\item Security Requirements and Verifiable Voting Schemes address different portions of the problem. A strong comparison states their scope, assumptions, enforcement points, evidence, and how a failure in one affects the other.
\item Reviewers should identify who controls voter, what inputs or dependencies it trusts, where privilege changes, what happens during partial failure, and how those assumptions are tested.
\item A weakness in voting can propagate across trust boundaries. The actual impact depends on reachable assets, granted authority, control independence, visibility, and recovery capability.
\item Useful evidence includes asset and dependency maps, physical-access records, sensor telemetry, safety constraints, maintenance logs, and recovery exercises; configuration alone is insufficient when operation, monitoring, or recovery is part of the claim.
\item Prevention reduces probability, detection limits dwell time, response limits spread, and recovery restores objectives. Omitting one layer creates an unexamined failure mode.
\item Begin by establishing scope, ownership, expected behavior, and the violated assumption. Then evaluate cyber and physical failure together, enforce safe states, and test operational recovery, preserving evidence for prioritization and follow-up.
\item Treat the exercise as stronger evidence than a static check. Prioritize according to reachable impact, likelihood of real operating conditions, missing compensating controls, and recovery difficulty.
\item Retain the principle, threat model, and control objective; separately label technologies, legal references, product examples, and threat observations whose applicability depends on time or environment.
\end{enumerate}
\section*{Key-Term Recap}
\begin{longtable}{>{\raggedright\arraybackslash}p{0.27\textwidth}>{\raggedright\arraybackslash}p{0.66\textwidth}}
\toprule\rowcolor{Navy}\color{white}\textbf{Term} & \color{white}\textbf{Meaning}\\\midrule
\endfirsthead
\toprule\rowcolor{Navy}\color{white}\textbf{Term} & \color{white}\textbf{Meaning}\\\midrule
\endhead
\textbf{Cipher} & An algorithm that transforms data using a key to provide a defined cryptographic security property. \\\midrule
\textbf{Encryption} & A keyed transformation of plaintext into ciphertext to protect confidentiality. \\\midrule
\textbf{Privacy} & The appropriate governance of information about people, including collection, use, disclosure, retention, and individual interests. \\\midrule
\textbf{Cryptography} & The use of mathematical mechanisms to protect confidentiality, integrity, authenticity, or related security properties. \\\midrule
\textbf{Integrity} & The property that information and systems remain accurate, complete, and protected from unauthorized modification. \\\midrule
\textbf{Threat} & A circumstance or actor capable of causing harm by exploiting a weakness or adverse condition. \\\midrule
\textbf{Digital Signature} & A public-key mechanism used to verify origin and integrity under defined key-trust assumptions. \\\midrule
\textbf{Availability} & The property that authorized users can access information and services when required. \\\midrule
\textbf{Confidentiality} & The property that information is not disclosed to unauthorized parties. \\\midrule
\textbf{Evidence} & Observable information that supports a security conclusion and can be preserved for review. \\\midrule
\bottomrule
\end{longtable}
\begin{CornellExamBox}{One-Sentence Takeaway}
Treat verifiable voting systems as a verifiable chain from assets and assumptions through controls, evidence, response, and recovery.
\end{CornellExamBox}
\end{document}
